\clearpage\thispagestyle{empty}
\pdfbookmark[0]{Quelle zum Bericht über Vesta, Juno, Pallas und Ceres}{planets1808july-sources}
\begin{center}\textbf{Quelle zum Bericht über Vesta, Juno, Pallas und Ceres}\end{center}
\noindent Diese KI-Übertragung ins heutige Deutsch folgt dem Bericht vom 4. Juli 1808 in \href{https://gdz.sub.uni-goettingen.de/id/PPN236007467}{Gauß’ \textit{Werken}, Bd.~VI} (Göttingen, 1874), S.~303--304. Der Bericht erschien unter der Überschrift \textit{Göttingen} im Doppelheft 107 und 108 der \textit{Göttingischen gelehrten Anzeigen}, \href{https://manifests.sub.uni-goettingen.de/iiif/image/gdz:PPN319721507_1808:00001071/full/1600,/0/default.jpg}{S.~1065}, \href{https://manifests.sub.uni-goettingen.de/iiif/image/gdz:PPN319721507_1808:00001072/full/1600,/0/default.jpg}{1066}, \href{https://manifests.sub.uni-goettingen.de/iiif/image/gdz:PPN319721507_1808:00001073/full/1600,/0/default.jpg}{1067} und \href{https://manifests.sub.uni-goettingen.de/iiif/image/gdz:PPN319721507_1808:00001074/full/1600,/0/default.jpg}{1068}. Die Werkausgabe nennt Stück 107. Der Bericht spricht von Gauß in der dritten Person.
\par\medskip
Im Erstdruck stehen auf S.~1067 zwei Beobachtungen der Juno. Die Werkausgabe verweist auf die entsprechenden Zeilen im Brief vom 27. Juni. Die Vestatafel des Berichts behält ihren Perihelwert, die eigene Genauigkeit der Angaben zu Knoten und Neigung und die unregelmäßigen Zeichen in den Punktreihen der Werkausgabe. KI-Übertragung und KI-Satz.
\clearpage
